\documentclass[aps,prb,onecolumn,nofootinbib,superscriptaddress]{revtex4-2}
\usepackage{amsmath,amssymb,amsthm,mathtools,bm}
\usepackage{booktabs}
\usepackage[colorlinks=true,linkcolor=blue,citecolor=blue,urlcolor=blue]{hyperref}
\usepackage{microtype}

\begin{document}

\title{CPU Choi Transfer-Gap Evidence in the Class-A Adaptive Fermionic Circuit}
\author{Worker 14}
\date{\today}

\begin{abstract}
This note summarizes read-only evidence for the CPU Choi/particle-transfer gap
analysis in the class-A adaptive fermionic circuit repository.  The strongest
CPU evidence is a production script and associated metadata for
domain-wall, slab-restricted Choi covariances, where the particle transfer
operator is recorded as
$T_p=-(I+\Sigma_{LL})^{-1}\Sigma_{LR}$ on finite sectors and the gap is
$\Delta_p=\min_j|\log s_j(T_p)|/C$.  The current CPU pipeline avoids saving full
Choi or transfer matrices but saves exact final spectra, near-gap eigenvectors,
all finite final eigenvectors in the newer campaigns, and spectrum-by-cycle
diagnostics.  The transfer-gap evidence is qualitatively consistent with the
independent Lyapunov analysis: both diagnostics identify near-marginal modes in
the domain-wall sector and motivate a chiral boundary interpretation.  The main
caveats are that this inspection did not load binary arrays, several numerical
claims therefore remain metadata-level, and some alpha-sweep runs show endpoint
saturation/NaN sentinel behavior that must not be mistaken for a resolved
finite gap.
\end{abstract}

\maketitle
\setcounter{tocdepth}{2}
\tableofcontents

\section{Source Scope}

This synthesis used only text files, JSON/CSV metadata, scripts, notebook text,
and existing notebook outputs.  No files were edited, no notebooks were
executed, and no large binary arrays such as \texttt{.npz} spectra or
eigenvectors were loaded.  Paths containing \texttt{haining}, \texttt{haoyu}, or
\texttt{erroneous} were excluded case-insensitively, as were
\texttt{repo\_synthesis/main\_pass/worker\_notes},
\texttt{supervisor\_notes}, and \texttt{boss}.

Primary inspected sources were \path{choi_covariance_cpu/},
\path{lyapunov_analysis_v2/majorana_choi_transfer_spectrum.ipynb}, the CPU Choi
campaign manifests and run summaries under
\path{choi_covariance_cpu/cpu_data/complex_particle_choi_transfer/campaigns/},
\path{scripts/smoke_choi_dense_vs_rank1.py}, and the existing Lyapunov writeup
\path{colab_lyapunov/docs/lyapunov_analysis.tex}.

\section{Findings: \texorpdfstring{$T_p$}{Tp} Gap and Choi Spectrum}

The CPU production entry point is explicit: the Choi transfer observable is
implemented in \path{choi_covariance_cpu/run_particle_choi_transfer_gap_cpu.py}
using the CPU class \texttt{classA\_U1FGTN}.  The script records
\texttt{classA\_U1FGTN.run\_markov\_circuit} as the canonical dynamics entry
point and defines the observable as a complex particle Choi-transfer gap.  Its
recorded transfer definition is
\begin{equation}
  T_p=-(I+\Sigma_{LL})^{-1}\Sigma_{LR},
\end{equation}
restricted to finite sectors.  The gap is recorded as
\begin{equation}
  \Delta_p(C)=\min_j\left|\frac{\log s_j(T_p)}{C}\right|.
\end{equation}
Operationally, the script diagonalizes the Hermitianized
$\Sigma_{LL}$ block, verifies that its spectrum lies in the pure-Choi interval
$[-1,1]$ up to tolerance, maps endpoint eigenvalues to particle zeros or poles,
and treats non-endpoint eigenvalues as finite particle-transfer exponents.

The relation between the Choi block eigenvalue $a$ and the finite transfer
exponent is implemented as
\begin{equation}
  \lambda_p(a;C)=\frac{\log(1-a)-\log(1+a)}{2C}.
\end{equation}
The near-gap modes are selected by increasing $|\lambda_p|$, with eigenvectors,
indices, residuals, endpoint multiplicities, finite-mode counts, and final
spectra stored by the observer.  The final-cycle gap is marked as a mandatory
exact eigensolve of $\Sigma_{LL}$ in the run summaries.

The most substantive CPU metadata campaign found here is
\path{Nx20_Ny20-30-40_DW1_dwtrunc1_a1-1-3_S10_cycles2Ny_allfiniteeig_v4}.  It
contains six domain-wall, slab-restricted, perfect-correction runs with
$N_x=20$, $N_y\in\{20,30,40\}$, $\alpha_1\in\{1,3\}$, $\alpha_2=30$, ten
samples per configuration, and $C=2N_y$.  The manifest and run index record all
six as completed, rank-1 physical covariance update used, no dense fallback, and
ten stable final samples.  The newer summaries also record that all finite final
eigenstates were saved and that full Choi/transfer matrices were not saved.

\section{Eigenvectors and Spatial Evidence}

The CPU observer saves near-gap final eigenvectors and, for newer campaigns, all
finite final eigenvectors.  The executed plotting notebook
\path{choi_covariance_cpu/plot_gap_vs_cycle_nan_sentinel_executed.ipynb}
constructs density maps by placing $|\psi|^2$ from active slab indices onto
$(x,y)$ and drawing domain-wall locator lines.  This is evidence that the saved
eigenvectors were intended for spatial localization diagnostics, but this note
does not independently inspect the binary eigenvector arrays.

The independent Lyapunov writeup gives the clearest text-level spatial
interpretation: the neutral Lyapunov eigenvector is reported to concentrate near
the two domain-wall interfaces and extend along the wall direction.  This is not
identical to the CPU particle-transfer eigenvector, but it supports the same
physical picture: a slow or marginal boundary-sector mode rather than a trivial
bulk instability.

\section{Relation to Lyapunov Evidence}

The Lyapunov analysis computes finite-time tangent exponents of the nonlinear
Markov map, whereas the CPU Choi analysis computes singular exponents of a
particle transfer object derived from the Choi covariance.  These are distinct
linearizations.  Nevertheless, both use a minimum absolute exponent as a gap
proxy.  The Lyapunov note states that DW-off spectra are strictly negative and
gapped, while DW-on spectra develop a near-zero exponent whose gap decreases
with $N_y$.  It also reports a sample-level correlation between small Lyapunov
gap and degraded Chern marker in the domain-wall sector.

The CPU Choi transfer-gap campaign is therefore best interpreted as a parallel
transfer-matrix diagnostic of the same slow domain-wall physics.  Agreement in
interpretation is qualitative at this read-only level: the metadata establishes
that exact Choi-block spectra and eigenvectors were produced and saved, while
the Lyapunov writeup supplies explicit text-level numerical trends and spatial
interpretation.

\section{Evidence Table}

\begin{table*}[t]
\caption{Text and metadata evidence inspected without loading binary arrays.}
\begin{tabular}{p{0.30\textwidth}p{0.19\textwidth}p{0.43\textwidth}}
\toprule
Source & Lines & Evidence used \\
\midrule
\path{choi_covariance_cpu/run_particle_choi_transfer_gap_cpu.py}
& 61--72
& Imports the CPU class, records the canonical CPU dynamics entry point, defines
the Choi formula, particle-transfer definition, gap definition, exact
$\Sigma_{LL}$ extraction, and rooted singular-value convention. \\
\path{choi_covariance_cpu/run_particle_choi_transfer_gap_cpu.py}
& 241--306
& Hermitianizes $\Sigma_{LL}$, diagonalizes it exactly, checks the $[-1,1]$
spectral interval, maps endpoints to zeros/poles, computes finite exponents,
selects near-gap modes, and records eigenvector residuals. \\
\path{choi_covariance_cpu/run_particle_choi_transfer_gap_cpu.py}
& 560--609
& Censors non-finite or out-of-bounds Choi blocks, performs exact final-cycle
storage of spectra and eigenvectors, and records all finite final eigenstates
when requested. \\
\path{choi_covariance_cpu/run_particle_choi_transfer_gap_cpu.py}
& 650--737
& Defines saved payloads: gap-vs-cycle, final spectrum, rooted singular values,
near-gap eigenstates, all finite eigenstates, endpoint counts, and stability
masks. \\
\path{choi_covariance_cpu/run_particle_choi_transfer_gap_cpu.py}
& 844--904
& Calls \texttt{classA\_U1FGTN.run\_markov\_circuit} with
\texttt{track\_choi=True}, perfect correction, slab measurement, Choi observer,
and rank-1/dense physical covariance update option. \\
\path{choi_covariance_cpu/.../Nx20_Ny20-30-40.../campaign_manifest.json}
& 2--17, 18--96
& Campaign metadata: $N_x=20$, $N_y=20,30,40$, $\alpha_1=1,3$,
$\alpha_2=30$, $C=2N_y$, ten samples, DW on, slab restricted, canonical CPU
entry point. \\
\path{choi_covariance_cpu/.../Nx20_Ny20-30-40.../campaign_manifest.json}
& 131--141, 147--165, 220--267
& Dense fallback policy and finite-mode policy are recorded; representative
results show completed runs with stable final samples and no rank-1 failure. \\
\path{choi_covariance_cpu/.../Nx20_Ny20-30-40.../run_index.csv}
& 1--7
& Six completed configurations are indexed; each has ten stable final samples,
rank-1 update used, and no dense fallback. \\
\path{choi_covariance_cpu/.../N20x40...a1-1.../run_summary.json}
& 888--893, 1127--1248, 1331, 1357
& Representative run summary: ten samples, censoring rule, exact final gap
override, saved spectrum-by-cycle and finite eigenstates, gap definition,
particle-transfer definition, physical covariance health, and transfer object
$T_pT_p^\dagger$. \\
\path{choi_covariance_cpu/.../N20x40...a1-3.../run_summary.json}
& 888--893, 1127--1248, 1331, 1357
& Same metadata for the $\alpha_1=3$ representative run; records ten stable
final samples and physical covariance eigenvalues essentially within
$[-1,1]$. \\
\path{choi_covariance_cpu/plot_gap_vs_cycle_nan_sentinel_executed.ipynb}
& 136--169, 823--837
& Existing executed notebook output loads six curves from the
$N_y=20,24,30$ campaign; $\alpha_1=1$ has no final sentinels, whereas
$\alpha_1=3$ has all samples sentinel by cycle 20 or 40, with first sentinel
cycles 11--13. \\
\path{choi_covariance_cpu/plot_gap_vs_cycle_nan_sentinel_executed.ipynb}
& 560--612
& Notebook constructs near-gap eigenstate density maps from saved eigenvectors
and active slab indices, with domain-wall locator lines. \\
\path{lyapunov_analysis_v2/majorana_choi_transfer_spectrum.ipynb}
& 7--10, 118--210, 233--260
& Older notebook converts complex Choi covariance to Majorana blocks and forms a
Majorana transfer spectrum using solve/pseudoinverse diagnostics; useful as
methodological context, less authoritative than the later CPU summaries. \\
\path{colab_lyapunov/docs/lyapunov_analysis.tex}
& 38--58, 206--221
& Defines Lyapunov gap proxy and reports DW-off large gaps versus DW-on
near-zero gaps decreasing with $N_y$; saves minimum-$|\lambda|$ eigenvectors. \\
\path{colab_lyapunov/docs/lyapunov_analysis.tex}
& 239--301, 464--471
& Reports strictly negative DW-off spectra, DW-on near-zero modes, slow
non-plateauing DW-on gap dynamics, and sample-level correlation between Chern
marker degradation and smaller Lyapunov gap. \\
\path{colab_lyapunov/docs/lyapunov_analysis.tex}
& 577--595, 625--651, 657--668
& Gives the chiral domain-wall interpretation, possible $1/N_y$ scaling, spatial
localization of the neutral mode near the domain walls, and open questions about
functional form and asymptotic convergence. \\
\path{scripts/smoke_choi_dense_vs_rank1.py}
& 1--6, 158--176, 191--215, 253--331
& Small-system health check comparing production rank-1 Choi resolvent action
against a dense solve/pseudoinverse implementation and reporting survival plus
final-block differences. \\
\bottomrule
\end{tabular}
\end{table*}

\section{Caveats}

The Choi transfer-gap arrays and eigenvector arrays were not loaded.  Thus this
note does not independently reproduce final gap values, rooted singular spectra,
or eigenvector density plots.  It reports the definitions and metadata-level run
status, plus text-level numerical conclusions already present in existing
documents or notebook outputs.

Endpoint saturation is a serious interpretive caveat.  The CPU policy accepts
cycles with fewer than the requested number of finite exponents and pads
near-gap slots with NaN or sentinel values rather than always censoring the
physical trajectory.  In the executed NaN-sentinel notebook, the $\alpha_1=3$,
$N_y=20,24,30$ curves have all samples at sentinel by cycle 20 or 40, so those
runs should be read as evidence for finite-sector collapse or numerical/endpoint
saturation, not as resolved finite transfer gaps.

The Majorana Choi transfer notebook predates the later CPU particle-transfer
pipeline and uses a pseudoinverse fallback when blocks are ill-conditioned.
This is useful for cross-checking ideas but is weaker evidence than the exact
$\Sigma_{LL}$ extraction in the current CPU script and run summaries.

The Lyapunov and Choi gaps are not the same spectrum.  The former is a tangent
spectrum of the nonlinear Markov map; the latter is a Choi-derived
single-particle transfer spectrum.  Their shared near-zero/domain-wall
interpretation is physically compelling but should be treated as convergence of
diagnostics, not equality of operators.

\section{Confidence}

Confidence is high for the implementation-level statements: the CPU script and
metadata explicitly define the transfer object, gap, exact extraction, saved
payloads, censoring policy, and campaign status.  Confidence is moderate for
the physical interpretation tying CPU Choi near-gap modes to chiral domain-wall
degrees of freedom, because that inference combines CPU metadata with the
separate Lyapunov writeup and notebook-derived plotting workflow.  Confidence is
low for any precise CPU Choi numerical scaling law in this read-only synthesis,
because the relevant binary spectra were intentionally not inspected.

\end{document}
